\documentclass{report}
\usepackage[utf8]{inputenc}
\usepackage{amsmath}
\usepackage{amssymb}
\usepackage{multicol}
\usepackage[a4paper, portrait, margin=1.2in]{geometry}
\usepackage{makeidx}
\makeindex

\title {Algebra Lineare e Geometria}
\author {Pietro Ballarin}

\begin{document}
\maketitle
\tableofcontents

\chapter{Introduzione}
Corso a cura di Mauro Costantini. Barbaro barbaroso

\section {Programma del corso}
\begin{itemize}
    \item Spazi vettoriali
    \item Sottospazi vettoriali
    \item Funzioni lineari
    \item Matrici
    \item Sistemi di equazioni lineari
    \item Autovalori ed autovettori
    \item Prodotti scalari
    \item Spazi affini ed euclidei
\end{itemize}

\chapter{Spazi vettoriali}

\section{Campi di numeri}

Lavoreremo con insiemi numerici dotati delle operazioni di somma ($+$)
e prodotto ($\cdot$) con le cosiddette \textbf{proprietà di campo}.
Per la somma:
\begin{itemize}
    \item Associativa
    \item Commutativa
    \item Esistenza elemento neutro ($0$)
    \item Esistenza dell'opposto ($-a$)
\end{itemize}
Per il prodotto:
\begin{itemize}
    \item Associativa
    \item Commutativa
    \item Elemento neutro ($1$)
    \item Distributiva
    \item Esistenza dell'inverso ($1/a$ oppure $a^{-1}$)
\end{itemize}
Un insieme con queste caratteristiche prende il nome di \textbf{campo}.
Baseremo la definizione di spazio vettoriale su questo concetto poiché
con queste proprietà siamo in grado di risolvere equazioni di primo grado,
essenziali per la risoluzione dei sistemi lineari e per la teoria stessa
della materia.

\section{Definizione}

Uno spazio vettoriale su un campo $K$ è un insieme non vuoto dotato di
due operazioni: \textbf{somma} ($+$) e \textbf{prodotto di un vettore per uno scalare} ($\cdot$).
Dove per \textbf{scalare} intendiamo un elemento del campo $K$ mentre per \textbf{vettore} intendiamo
un elemento dello spazio vettoriale.

\section{Esempi di spazi vettoriali}
Alcuni esempi presenti negli appunti del docente:
\begin{itemize}
    \item{Le n-uple di elementi di un campo $K$}
    \item{I polinomi}
\end{itemize}

\section{Combinazioni lineari}
Sia $V$ uno spazio vettoriale sul campo $K$ siano $v_1, ..., v_n \in V$ e $\lambda_1, ..., \lambda_n \in K$
si dice combinazione lineare dei vettori di $V$ il vettore:
$$ \lambda_1 v_1 + \lambda_2 v_2 + ... + \lambda_n v_n $$

\section {Vettori linearmente dipendenti e indipendenti}

\subsection {Vettori linearmente indipendenti}
I vettori $v_1, ..., v_n$ si dicono linearmente indipendenti se l'unica soluzione di
$ \lambda_1 v_1 + \lambda_2 v_2 + ... + \lambda_n v_n = 0 $ e' valida solo per
$\lambda_1 = \lambda_2 = ... = \lambda_n = 0$

\subsection {Vettori linearmente dipendenti}
I vettori si dicono linearmente dipendenti se non sono linearmente indipendenti ovvero se esistono
dei coefficienti non tutti nulli tali che la loro combinazione lineare e' nulla

\end{document}
